\documentclass{article}
\usepackage{/globals/anki}
\ankiprefix{pk-systeme}
\ankitopic{Planungs- und Kontrollsysteme}
\ankishow
\begin{document}
\section{Planungs- und Kontrollsysteme} 
Planungs- und Kontrollsysteme können genutzt werden um Engpässe und Überkapazitäten zu bestimmen, um Kapazitätsbedarfe zu berechnen.

\subsection{Lebenszyklus}
IT-Produkte haben generelle kürzere Lebenszyklen. In Bezug auf den Lebenszyklus muss geplant werden ob ein Produkt die einen Entwicklungskosten decken kann, die Entwicklungskosten des Nachfolgeproduktes decken kann und die Wartungskosten decken kann.
\ankinote{lebenszyklus}{lebenszyklus}{Eigene Kosten, Nachfolgerkosten, Wartungskosten}

\subsection{Absatzzyklen}
Abhängig von der Art des Produktes ist es zu erwarten dass der Absatz in Abhängigkeit von von der Zeit anders aussieht.
\begin{description}
 \item[Hardware] Neue Hardware verkauft sich über einen längeren Zeitraum gut, oftmals so lange bis es einen Nachfolger gibt. Dann sinkt die Absatzmenge schnell.
 \item[Hype] Hypeprodukte verkaufen sich anfangs massiv, werden aber schnell irrelevant.
 \item[Langsame Produktion] Manche Produkte, vor allem neue Produkte, brauchen eine Weile bis die Produktionsgeschwindigkeit ihr Maximum erreicht.
 \item[Saisonale Produkte] Eine Sinuskurve.
 \item[Neue Produktarten] Solange die Produktion und das öffentliche Intresse steigt steigt auch der Absatz exponentiell.
\end{description}
\ankinote{absatz}{Absatzzyklen}{Hardware, Hype, Langsame Produktion, Saison, Neue Produktart}

\end{document}